% !TEX root = ../main.tex
\section{Experimental Setup}
\label{sec:experiments}

This chapter describes how the nine models of Chapter~\ref{section:methods} are
trained and evaluated. It specifies the tools and hardware, the training
configuration and the experiment set, and the held-out validation protocol. It
then describes the cross-environment comparison against a second simulator, an
off-policy baseline based on Soft Actor-Critic, and the metrics reported in the
results.

\subsection{Tools and Infrastructure}
\label{sec:tools}

Training and Isaac-Lab evaluation run on a single RTX Pro 6000 GPU with 96\,GB of
memory. The software stack is Isaac Sim 5.1.0, Isaac Lab 2.3.0, cuRobo 0.7.5 for
batched Inverse Kinematics (IK), and PyTorch 2.7.0. A second environment, the
two-dimensional \texttt{gym-pusht} simulator, is used for the cross-environment
comparison of Section~\ref{sec:cross_env}. It runs on 32 CPU cores and uses
neither IK nor three-dimensional contact physics. The off-policy baseline of
Section~\ref{sec:sac_baseline} uses Stable-Baselines3 2.7.0. Running the same
models in two simulators of different fidelity separates conclusions that hold
across environments from artefacts of a single implementation.

\subsection{Training Configuration and the Experiment Set}
\label{sec:train_config}

At the 528-environment anchor, every model is trained for approximately 3{,}000
iterations. This corresponds to about 1.6 days of wall-clock time per model on the
RTX Pro 6000. At 15 pushes per rollout it yields a batch of roughly 7{,}920 push
transitions per PPO update. For the reward-efficiency comparison, the
improvement-reward baseline and the potential-based agent are trained across the
same environment counts, from 256 to 2{,}048, so the two rewards are measured at
matched parallelism rather than at a single batch. The two self-play subjects are
additionally trained at 256 and 2{,}048 environments, which tests whether the
single-GPU compute budget is the cause of their result. The PPO hyperparameters are
a learning rate of $3\times10^{-4}$, a clipping range of 0.2, a discount factor of
0.998, a GAE parameter of 0.95, and an entropy coefficient of 0.005. The self-play
models additionally use Alice Behavioral Cloning at $\beta = 0.5$.

Two points about the training history are stated for reproducibility. First, all
reported checkpoints post-date a sequence of corrections to the reward and the
environment. The reported numbers are therefore produced on the corrected code,
and earlier runs are excluded. Second, the improvement-reward baseline and the
potential-based agent are compared across a shared range of environment counts, so
that any difference between the two rewards cannot be attributed to unequal
parallelism. The budgets are tabulated in Appendix~K.

Every Isaac Lab run in the experiment set is budget-matched. The total experience
is held at about 23.8 million pushes, and the iteration count is set inversely
proportional to the number of parallel environments, so a change in the number of
environments changes the per-update batch size alone. This design has one confound that is
stated wherever a scale result appears: a larger batch also means fewer gradient
updates over training, so batch size and update count move together and cannot be
fully separated. Table~\ref{tab:phases} lists the phases, the models each one
trains, the environment counts, the matched iteration counts, and the seeds.

\begin{table}[H]
    \centering
    \small
    \begin{tabularx}{\textwidth}{l X c c c c}
        \hline
        \textbf{Phase} & \textbf{Models} & \textbf{Envs} & \textbf{Iterations} &
        \textbf{Seeds} & \textbf{Runs} \\
        \hline
        Anchor and efficiency & PPO-Baseline, PPO-PBRS & 256--2{,}048 &
        773--6{,}188 & 3 & 24 \\
        Seed intervals & Baseline, PBRS, Curriculum, ASP-dPose, TASP-dPose & 528 &
        3{,}000 & 5 & 25 \\
        Coupling isolation & ASP-disc, TASP-disc & 528 & 3{,}000 & 5 & 10 \\
        Self-play scale & ASP-dPose, TASP-dPose & 256, 2{,}048 & 6{,}188; 773 & 3 &
        12 \\
        Component ablations & ASP-dPose, TASP-dPose & 528 & 3{,}000 & 3 & 12 \\
        Imitation positive control & ASP-disc & 528 & 3{,}000 & 3 & 6 \\
        Reward structure & Bob-penalty & 528 & 3{,}000 & 5 & 5 \\
        Cross-environment & single agent, curriculum, self-play & 64, 256 &
        3{,}000 & 1 & 10 \\
        \hline
    \end{tabularx}
    \caption{The phases of the experiment set. Environments and iterations trade
    off to hold the total experience at about 23.8 million pushes for the Isaac Lab
    runs; the cross-environment runs instead use a fixed 3{,}000 iterations and a
    single seed, and vary the per-update batch through the environment count and
    the push length. The experiment set comprises 94 training runs in Isaac Lab and
    10 in the second simulator.}
    \label{tab:phases}
\end{table}

\subsection{Validation Protocol}
\label{sec:val_protocol}

Every model is evaluated on the same held-out suite of 30 T-block scenes,
illustrated in Figure~\ref{fig:val_layout}. The suite is divided into three groups
of ten. The first ten are rotation-dominant scenes with large orientation
changes. The next ten are position-only scenes whose goal orientation equals the
start orientation. The last ten are combined scenes that require both a position
and an orientation change. Within each group the difficulty ranges from easy to
hard, graded by the start-to-goal distance and the magnitude of the required
rotation.

\begin{figure}[H]
    \centering
    \includegraphics[width=\textwidth,height=0.42\textheight,keepaspectratio]{val_test_layout.png}
    \caption{The 30-scene held-out validation suite. Each panel shows the block
    start pose and the goal pose for one scene. The suite is identical for every
    model.}
    \label{fig:val_layout}
\end{figure}

Each scene is attempted 20 times, and within each attempt the policy is allowed up
to 30 pushes. Both the single-agent and the self-play validators sample actions
from the policy, rather than taking the most probable action, so every model is
scored under one protocol. Two success measures are reported and kept distinct
throughout. The \emph{scene success rate} counts a scene as solved if at least one
of its 20 attempts satisfies the combined gate of Section~\ref{sec:mdp}. This is
the headline number, and follows the convention common in manipulation evaluation.
The \emph{trial success rate} is the fraction of individual attempts that satisfy
the gate, and is always lower than the scene success rate. Reporting both prevents
the best-of-20 protocol from overstating performance, and the gap between them is
stated wherever the headline number appears; the final values are \TBD until the
seed runs are collated. Each model is evaluated across its trained seeds, and the
results are aggregated into a mean with a 95\,\% confidence interval per model,
batch size, and iteration count.

\subsection{Cross-Environment Protocol}
\label{sec:cross_env}

To test whether the ordering of the models is a property of the methods rather
than of Isaac Lab, the single-agent, curriculum, and self-play models are
retrained and evaluated in \texttt{gym-pusht} under the same 30-scene gate. The
two environments differ mainly in the size of the PPO batch. Isaac Lab produces
about 7{,}920 transitions per update at the 528-environment anchor, while
\texttt{gym-pusht} produces between about 960 and 11{,}520, depending on the
environment count and the push length. The gym runs use a fixed 3{,}000 iterations
and sweep the per-update batch through these two settings, so the comparison tests
whether the ordering of the methods survives a second, lower-fidelity simulator,
rather than matching the total experience between the two environments.
Figure~\ref{fig:isaac_vs_gym} summarises the throughput and batch size of the two
environments.

\begin{figure}[H]
    \centering
    \includegraphics[width=0.7\textwidth,height=0.4\textheight,keepaspectratio]{isaac_vs_gym.png}
    \caption{Throughput and PPO batch size for Isaac Lab and \texttt{gym-pusht}.
    Raw throughput is comparable; the per-update batch differs by roughly a factor
    of eight, which is the variable the cross-environment comparison isolates.}
    \label{fig:isaac_vs_gym}
\end{figure}

\subsection{Off-Policy Baseline: SAC with Hindsight Experience Replay}
\label{sec:sac_baseline}

The principal models are all on-policy. To provide an external point of
calibration from a different algorithm family, an off-policy baseline is built
from Soft Actor-Critic (SAC) \cite{Haarnoja2018SoftAO} combined with Hindsight
Experience Replay (HER) \cite{Andrychowicz2017}. SAC is a natural counterpart to
PPO for this problem. It is sample-efficient through experience replay, and its
maximum-entropy objective (Equation~\ref{eq:sac}) supports exploration under
sparse reward. HER addresses the sparse goal-conditioned setting directly: it
relabels the object pose actually reached in a failed attempt as if it had been
the goal, so every attempt yields usable transitions.

The baseline runs on a direct control environment through Stable-Baselines3, and
reuses the shared push primitive and the cuRobo IK pipeline of
Section~\ref{sec:mdp}. Two adaptations are required. First, SAC is a
continuous-action algorithm, so the four push parameters become a continuous
action in $[-1, 1]^4$ and are decoded to the same object-relative primitive,
rather than the MultiCategorical variable used by the PPO models. Second, the
goal-conditioned observation is expressed in the dictionary form the replay buffer
expects, with separate observation, achieved-goal, and desired-goal entries. The
baseline is trained and evaluated under the same 30-scene protocol as the
principal models, so its success rate is directly comparable. Its role is
calibration rather than competition. It indicates whether an off-policy,
replay-based method reaches performance similar to the on-policy potential-based
model, and its status against the compute budget is reported with the results.

\subsection{Metrics}
\label{sec:metrics}

The results report the following quantities, each with its sample size and
protocol stated at the point of use. The primary metric is the scene success rate
over the 30-scene suite, with the trial success rate reported alongside it. The
breakdown by scene type, position-only against combined, is reported because it
localises where each model fails. Mean position error and mean orientation error
are reported as continuous measures of near-misses. For the self-play models the
training logs also provide the per-axis success rates: the fraction of pushes that
satisfy the position gate alone and the orientation gate alone. These are compared
against the combined-gate rate to diagnose whether the two skills are learned but
not combined. Average push count per attempt is reported as a measure of
efficiency. Three further quantities are logged during self-play training and
reported in the results. The first is the mean absolute value error, which
measures the bias of the value function under the shifting goal distribution. The
second is the share of demonstrated actions the imitation term clips away, a direct
measure of how strongly the imitation signal is suppressed. The third is the
Alice-replay rate: the fraction of proposed goals that the goal-proposing agent's
own action sequence reaches when replayed from the goal-solving agent's start
state. All reported numbers trace to the validation records summarised in the
results.
